\section*{Resumen}

\noindent Mientras la sociedad permanece atrapada en el ciclo del consumismo, del que
parece no haber salida, y las empresas continúan fomentándolo, surge una
pregunta: ¿cómo puede una aplicación móvil de escaneo de códigos de barras
contribuir a reducir la asimetría de información en la compra de productos
electrónicos y fomentar un consumo más informado y sostenible? Responderla es
importante para entender cómo se comporta la sociedad cuando dispone de
herramientas que le permiten tomar decisiones informadas y que, además, reducen
el sesgo informativo generado por las campañas de marketing de las empresas. La
creación de la aplicación ``TreeBar'' se respalda en una encuesta aplicada a más
de 200 personas mayores de edad, cuyo objetivo no fue solo evaluar la viabilidad
de la aplicación, sino también comprender cómo se comporta la sociedad actual en
el mundo de las compras. Esto es relevante porque conocer cómo toman sus
decisiones los compradores permite diseñar correctamente una herramienta que
busca transformar los hábitos de consumo de los usuarios de laptops y
smartphones. Los resultados fueron valiosos para el desarrollo del proyecto: el
80\% de los encuestados manifestó interés en una aplicación de esta naturaleza,
la cual no solo muestra información concreta de cada producto, sino que además
permite consultar reseñas de usuarios reales, formular preguntas para que las
respondan quienes ya adquirieron el dispositivo y, finalmente, encontrar lugares
donde repararlo.

\noindent\textbf{Palabras clave:} Sociedad, empresas, consumismo, informado, concreta.